\section{Introduction}
\label{sec:intro}

Lens flare \cite{fowles1989introduction,hullin2011physically,wu2021train}, an optical artifact caused by reflections and scattering within a camera lens, is a well-known source of image degradation in conventional photography \cite{wu2021train,zhou2023improving,dai2022flare7k,dai2024flare7k++,jiang2024mfdnet}. It severely compromises downstream vision tasks and is not fully avoidable by hardware such as lens hoods, which \textbf{cannot block light sources within the field of view}, motivating extensive and ongoing research on lens flare removal in the image domain. 


Event cameras \cite{gallego2020event,li2021recent,posch2014retinomorphic, chaney2023m3ed,perot2020mpx1,kim2021n-imagenet,muglikar2025event}, which asynchronously capture intensity changes at microsecond resolution, are increasingly important for high-speed and high-dynamic-range applications \cite{tulyakov2022time,hagenaars2021self,li2022asynchronous,lin2024e2pnet,vemprala2021representation,wang2023visevent,zheng2023deep,gehrig2024low,han2024event,li2025active,alonso2019ev-segnet,sun2022ess,kong2024openess,kong2025eventfly,kong2025talk2event,lu2025flexevent,sun2024unified,gehrig2024dense,brander2025reading,messikommer2025data,geckeler2025event}. However, their high dynamic range does not extend to seeing through lens flare, whose intensity is superimposed onto the scene irradiance before the sensor, \textbf{rendering the two inseparable} \cite{brandli2014davis,posch2010gen1,son2017gen3,finateu2020gen4}. Flare-like artifacts are present in many widely used event datasets~\cite{gehrig2021dsec,gehrig2021raft,hu2020ddd20,binas2017ddd17,zhu2018multivehicle}, where they degrade data quality and hinder downstream performance~\cite{shi2023identifying,shi2024polarity,hidalgo2020learning,zubic2023from,jing2024hpl-ess,hamaguchi2023hmnet}. These artifacts have been attributed to generic lighting interference~\cite{shi2023identifying}, without recognizing them as lens flare. We argue that lens flare in event cameras is \textbf{a critical yet overlooked problem}. As shown in Fig.~\ref{fig:teaser}, it manifests as a highly structured, spatio-temporal distortion precisely explained by the underlying optics.


This oversight has significant consequences. In event streams, lens flare appears as large-scale, temporally correlated patterns that differ fundamentally from \emph{random noise} or \emph{periodic flickers}. Existing event-based denoising methods \cite{shi2024polarity,shi2023identifying,wang2022linear,guo2022low,duan2021eventzoom}, designed for \emph{uncorrelated} or \emph{uniform noise}, fail to suppress structured artifacts. Moreover, strong flare signals can \textbf{suppress or distort valid background events} \cite{valencia2023optically}. Recovering this corrupted background is akin to event-based see-through \cite{yu2022learning}, yet far more challenging: those methods assume controlled settings in which only the occluder moves, whereas lens flare is an inherently unconstrained corruption present in arbitrary scenes.


A learning-based approach appears most promising, yet it is hindered by a fundamental obstacle: \textbf{the absence of paired training and testing data}~\cite{dai2022flare7k,dai2024flare7k++}. Capturing perfectly aligned event pairs with and without flare requires re-recording the same scene under identical conditions, which is infeasible. This challenge is amplified for event cameras, whose motion encoding makes trajectory alignment nearly impossible. Naively applying an event simulator~\cite{hu2021v2e,lin2022dvs,joubert2021event,zhang2023v2ce,han2024physical} to existing image datasets yields a fully synthetic domain for both input and ground truth, creating a severe sim-to-real gap. Simulation-based approaches in the image domain have addressed this scarcity by collecting flare and background assets \emph{separately}, then compositing them into paired training samples via direct linear superposition~\cite{wu2021train,dai2022flare7k,dai2024flare7k++}, a strategy whose practicality is now well established. The analogous route for event cameras is natural, yet its compositing step is categorically different. Event cameras encode irradiance in the logarithmic domain, merging background and flare event streams is \textbf{not a simple union but a non-linear suppression process} that has never been formally analyzed or characterized.



We introduce \emph{\ours}, a framework that resolves this challenge by grounding the data generation pipeline in physical analysis. We derive a physical forward model revealing the non-linear compositing mechanism at the core of this problem, which directly guides our primary contributions: a physics-driven simulator that generates the large-scale \textbf{E-Flare2.7K} training set, and a novel, physically-principled acquisition system that captures \textbf{E-Flare-R}, the first-ever paired real-world test set. This benchmark enables training \textbf{E-DeflareNet}, a standard 3D U-Net restoration network, to restore clean event streams with high fidelity.


Our main contributions can be summarized as follows:
\begin{itemize}
    \item To our knowledge, we are the first to provide a rigorous physical analysis of lens flare in event data, culminating in a \textbf{forward model} that formalizes the non-linear suppression mechanism underlying event stream corruption. This builds a theoretical foundation for de-flaring tasks in event-based vision.

    \item Guided by our physical model, we introduce the first \textbf{simulation pipeline} that explicitly models non-linear flare corruption in event streams, enabling large-scale, physically-grounded paired training data for the first time.

    \item We design and contribute the \textbf{E-Deflare Benchmark}, the first comprehensive resource for the event camera de-flaring task, comprising a large-scale simulated dataset (\emph{E-Flare2.7K}), a curated set of real-world examples (\emph{DSEC-Flare}), and critically, the first-ever paired \textbf{real-world test set} (\emph{E-Flare-R}) for comprehensive evaluations and assessments.

    \item We use a standard 3D U-Net, \textbf{E-DeflareNet}, to validate the proposed physical formulation and paired data, showing that training on our data yields strong de-flaring performance and improves downstream event camera imaging.
\end{itemize}



















% Lens flare \cite{fowles1989introduction,hullin2011physically,wu2021train}, an optical artifact from internal scattering and reflections within a camera's lens, is a well-known source of degradation in conventional photography~\cite{wu2021train,zhou2023improving,dai2022flare7k,dai2024flare7k++,jiang2024mfdnet}. It severely compromises downstream computer vision tasks, motivating substantial research on flare removal in the image domain. Event cameras~\cite{gallego2020event, li2021recent, posch2014retinomorphic}, which asynchronously record light changes at the microsecond level, are increasingly critical in challenging vision tasks~\cite{tulyakov2022time,hagenaars2021self, li2022asynchronous, lin2024e2pnet, vemprala2021representation, wang2023visevent, zheng2023deep,gehrig2024low,han2024event,li2025active}. While their unique sensing paradigm offers advantages like high dynamic range, they are not immune to this fundamental optical limitation. Indeed, flare-like artifacts are visibly present in many prominent event camera datasets~\cite{gehrig2021dsec,gehrig2021raft,hu2020ddd20,binas2017ddd17,zhu2018multivehicle}, where they degrade data quality and affect downstream tasks~\cite{shi2023identifying,shi2024polarity,hidalgo2020learning}. However, these artifacts have generally been attributed to generic light source interference~\cite{shi2023identifying}, without a systematic analysis that identifies them as the distinct optical phenomenon of physical lens flare. We argue that lens flare in event cameras is a critical and overlooked problem. As illustrated in Fig.~\ref{fig:teaser}, it manifests as a highly complex spatio-temporal phenomenon, whose intricate patterns can be precisely explained by the underlying physics of lens flare.

% This oversight has significant consequences. In event streams, lens flare manifests as large-scale, complex, and \textbf{spatio-temporally correlated} patterns. Existing event-based denoising methods~\cite{shi2024polarity, shi2023identifying, wang2022linear,guo2022low,duan2021eventzoom}, designed for random, uncorrelated noise or uniform flicker, are fundamentally ill-equipped to handle such structured artifacts. The contamination is not limited to adding spurious activity; as our theoretical analysis will explain, the strong flare signal can also suppress and distort legitimate event data from the background~\cite{muglikar2025event,valencia2023optically}, making the de-flaring task even more difficult.

% To tackle such a complex artifact, a learning-based approach is the most promising path. However, this path is blocked by a fundamental obstacle: the absence of suitable training and testing data~\cite{dai2022flare7k,dai2024flare7k++}. Generating large-scale, perfectly-aligned paired data would require capturing a dynamic scene twice with pixel-perfect alignment—once with flare and once without. This is an extremely difficult task, as one cannot simply disable the optical effect without altering the scene's illumination, nor can one perfectly replicate the precise trajectory of any moving objects. This data acquisition challenge is even more severe for event cameras, which inherently capture dynamic information~\cite{gallego2020event}, making the alignment of trajectories exceptionally difficult. Furthermore, a naive workaround, such as applying an event simulator~\cite{hu2021v2e,lin2022dvs,joubert2021event,zhang2023v2ce,han2024physical} to existing paired frame datasets, is also untenable, as it would create a fully simulated dataset for both input and ground truth, introducing a significant sim-to-real gap.

% To break this data deadlock, this paper presents \textbf{E-Deflare}, a systematic framework built upon a rigorous, physics-driven methodology. Our approach is founded on a deep analysis of the underlying optical principles, which enables a complete solution from theory to practice. We first establish a theoretical foundation for the problem, then leverage it to engineer a data-centric solution that makes learning-based de-flaring possible. We demonstrate the success of this framework by training a 3D U-Net-based architecture~\cite{10.7554/eLife.57613,cciccek20163d,lee2017superhuman} to effectively restore clean event streams.

% Our main contributions are as follows:
% \begin{itemize}
%     \item \textbf{Theoretical Foundation.} We are the first to systematically identify lens flare in event cameras and derive a theoretical \textbf{forward model} that explains how it contaminates event streams via a non-linear suppression mechanism. This provides the physical basis for designing effective de-flaring strategies.
    
%     \item \textbf{Simulation for Training.} Based on our model, we build the first physics-driven \textbf{simulation pipeline} to generate a large-scale, paired dataset. This overcomes the primary bottleneck for training deep de-flaring networks.
    
%     \item \textbf{Real-World Validation and Benchmark.} We propose a novel method to acquire a paired \textbf{real-world test set}, enabling robust validation of sim-to-real generalization. We will release our synthetic and real datasets as the first public \textbf{benchmark} to foster future research in this area.
% \end{itemize}
